\documentclass[a4paper, 12pt]{article}
\usepackage{graphicx}

\usepackage[utf8]{inputenc}
\usepackage[magyar]{babel}

% a szép matematikai szimbólumokért
\usepackage{amssymb}
\usepackage{amsmath}

% ha táblázatban szeretnénk egyesített sorokat is
\usepackage{multirow}

% horizontal line
\usepackage{hhline}

% eps formátumú ábrák --> pdflatex fordításhoz!!
\usepackage{epsfig}

% egyenletekhez, pl mátrixok írására
\usepackage{array}

% ha betűszíneket is szeretnénk használni
\usepackage{color}

% Margók egyéni beállításai
\usepackage{anysize}
\marginsize{1.64cm}{1.64cm}{1.2cm}{2.4cm} %\left right top bottom

% vakszöveg
\usepackage{lipsum}
\usepackage{blindtext}

% A HIVATKOZÁSOKHOZ HASZNÁLT CSOMAGOK. RÉSZLETESEBBEN LD: google -> latex bibtex
\usepackage[numbers, square, comma, sort&compress]{natbib}
%\usepackage[format=hang,labelsep=period]{caption}

\usepackage[unicode]{hyperref}   % ezzel a hivatkozások linkké válnak
\usepackage{bookmark}
\hypersetup{bookmarksopen={true}}
\hypersetup{bookmarksopenlevel={2}}
\hypersetup{bookmarksnumbered={true}}
\hypersetup{
  colorlinks,%
  citecolor=red,%
  filecolor=black,%                                                                                                                                               
  linkcolor=blue,%
  urlcolor=green
}
%\numberwithin{equation}{section}          % ezekkel tudod beállítani, hogy milyen felbontásig menjen a hivatkozás
%\numberwithin{figure}{subsection}
%\numberwithin{table}{section}          % ha kikommenteled, akkor csak simán számozva lesz.

%%%%%%%%%%%%%%%% Néhány dolog a fancy kinézethez

\frenchspacing
\setlength{\parskip}{2ex}
\setlength{\headsep}{0,4cm}
\setlength{\headheight}{4pt}

% fej- es lábléc
\usepackage{fancyhdr}
\usepackage{fancyref}
\usepackage{fancyvrb}
\pagestyle{fancy}

\renewcommand{\headrulewidth}{0,05pt}
\renewcommand{\footrulewidth}{0pt}


\fancyhf{}
\fancyhead[RE]{{ \nouppercase{\leftmark}} }
\fancyhead[LO]{{ \nouppercase{\leftmark}} }
\cfoot{--~\thepage~--}


%%%%%%%%%%%%%%%%%%%%%%%%%%%


%%%%%%%%%%%%%%%%%%%%%%%%%%%%%%%%%%%%%%%%%%%%%%%%%%%%%%%%%%%%%%%%%%%%%%%%%%%%

\begin{document}

% Címoldalt lehet egyszerűen a \maketitle paranccsal is. Ha kissé részletesebb
% címre van szükség, azt lehet így is, kézzel megadva mindent.
\begin{titlepage}   
\begin{center}
\thispagestyle{empty}  

\vspace*{0.7cm}
\rule{\linewidth}{0.5mm} \\[3mm]
\vspace*{0.7cm}

{\LARGE Kutatómunka információs eszközei}

\vspace*{0.7cm}
\rule{\linewidth}{0.5mm} \\[3mm]
\rule{\linewidth}{0.5mm} \\[3mm]



{\Large Differenciálegyenletek numerikus integrálási módszerei\\}

\vspace*{0.7cm}
\rule{\linewidth}{0.5mm} \\[3mm]
{\footnotesize Írta: Rácz Gergely} \\
{\tiny graczgeri@gmail.com}
  \vspace*{2cm}

\begin{figure}[h!]
\begin{center}
\includegraphics[width=0.5\textwidth]{img/elte}
\end{center}
\end{figure}

\end{center}
\end{titlepage}

\newpage

%%%%%%%%%%%%%%%%%%%%%%%%%%%%%%%%%%%%%%%%%%%%%%%%%%%%%%%%%%%%%%%%%%%%%%%%

\thispagestyle{empty}  
  \vspace*{3cm}
\begin{abstract}
A megírt programok a Föld-Hold rendszer mozgásegyenletét oldják meg, numerikus integrálási módszerekkel. Két módszert, az explicit Euler módszert és az  adaptív lépéshossz-szabályozott negyedrendű Runge-Kutta-módszert hasonlítottam össze.
\end{abstract}

% ---------------------------- T A R T A L O M J E G Y Z É K ----------------------------
\newpage \vspace*{1cm}
\thispagestyle{plain}                                                                                                                                             
%\pagenumbering{roman} \setcounter{page}{1}
\tableofcontents

\vspace*{1cm}
\listoffigures

\vspace*{1cm}
%\listoftables

\newpage
\pagestyle{fancy}
\pagenumbering{arabic} \setcounter{page}{1}
\vspace*{1cm}

% ---------------------------- M A G A   A   D O L G O Z A T ----------------------------

%%%%%%%%%%%%%%%%%%%%%%%%%%%%%%%%%%%%%%%%%%%%%%%%%%%%%%%%%%%%%%%%%%%

\section{Bevezetés}
\label{sec:bev}

A Föld-Hold rendszerének mozgását vizsgáltam síkban. A feladat során geocentrikus koordinátákkal dolgoztam. A mozgást Newton általános mozgástörvénye írja le:
\begin{center}
\begin{equation}
\label{newtequ}
F=\gamma\frac{m\cdot M}{r^2}
\end{equation}
\end{center}
Az Euler módszer során a Földet rögzítettem az origóban (0;0), a negyedrendű Runge-Kutta módszer során a Föld és a Hold is szabadon mozoghatott, a megadott kezdőfeltételeknek megfelelően. A feladat tulajdonképpen differenciálegyenletek numerikus integrálása, különböző pontosságú módszerekkel.



%%%%%%%%%%%%%%%%%%%%%%%%%%%%%%%%%%%%%%%%%%%%%%%%%%%%%%%%%%%%%%%%%%%%%%%%%%%%%%%%%%%%%%%
\section{Program}
\label{sec:prog}

A programokat C nyelven írtam, először CodeBlocks környezetben a \cite{artic:negyedik} forrás alapján, majd Visual Studio Code-ban pontosítva, alakítva a jelenlegi feladathoz. A program írása közben Git-es verziókövetést használtam, a repositoryt a \url{https://github.com/graczgeri/foldhold} link alatt lehet elérni. A program írása közben több branchet is létrehoztam, az áttekinthetőség kedvéért, ezeket - miután nem volt rá szükség és illesztettem a master brachbe a megfelelő részeket - töröltem. Az összehasonlítás szemléltetésére gnuplot segítségével plotoltam a kapott adatsorokat. Az egészet CMake segítségével buildeltem. A használat során a \cite{artic:elso,artic:masodik} forrásokra támaszkodtam. A CMake használata végett külön scriptet írtam az ábrák elkészítésére és egy előre gyártott tex mintát módosítottam a saját felhasználásnak megfelelően. Sajnos a pdf file elkészítése nem sikerült a CMake használatával, így azt külön csináltam meg.

\subsection{Euler módszer}
\label{subsec:euler_1}

Az Euler módszer lényege, hogy $\delta$x (általában h-val jelölt) diszkrét lépéshosszal léptetjük a változók értékét. Az összes változót léptetjük egyszerre:
\begin{center}
\begin{equation}
y_{n+1}=y_n+h\cdot f(x_n,y_n)
\end{equation}
\end{center}
Eközben a független változót is léptetjük:
\begin{center}
\begin{equation}
x_{n+1}=x_n+h
\end{equation}
\end{center}
A módszer hibája, hogy ha a derivált a lépés során túl gyorsan változik, akkor a lépésnek relatív nagy lesz a hibája. Sok lépés után nagy lesz az eltérés az analitikus megoldástól, mivel a hiba a lépések számával felösszegződik.

\subsection{Runge-Kutta 4 módszer}
\label{subsec:rk4_1}

Az Euler módszer javításának tekinthető a középponti módszer. Maga az Euler módszer aszimmetrikus, mivel a deriváltat mindig az $x_n$ helyen vesszük. A középponti módszer lényege, hogy teszünk egy fél lépést, kiszámoljuk a deriváltat a középső pontban és ezt használjuk a teljes lépés során. A Runge-Kutta módszer ezt használja fel, több lépésben:
\begin{center}
\begin{equation}
k_1=h\cdot f(x_n,y_n)\newline
\end{equation}
\begin{equation}
k_2=h\cdot f(x_n+\frac{1}{2}h,y_n+\frac{1}{2}k_1)\newline
\end{equation}
\begin{equation}
k_3=h\cdot f(x_n+\frac{1}{2}h,y_n+\frac{1}{2}k_2)\newline
\end{equation}
\begin{equation}
k_4=h\cdot f(x_n+h,y_n+k_3)\newline
\end{equation}
\begin{equation}
y_{n+1}=y_n+\frac{1}{6}k_1+\frac{1}{3}k_2+\frac{1}{3}k_3+\frac{1}{6}k_4+O(h^5)
\end{equation}
\end{center}
Amiben $k_{(1..4)}$ a lépések, $O(h^5)$ pedig a lecsökkent hiba az analitikus eredménytől. A lépések együtthatóit a Butcher-tábla adja.

\section{Eredmény}
\label{sec:ered}

\subsection{Euler módszer}
\label{sebsec:euler_2}

A kiértékelés során a program az euler.dat adattömbbe írta a Hold x és y koordinátáit. A kapott adatokat ábrázolva látható az \ref{fig:euler}. ábrán a Hold Föld körüli pályája, rögzített Föld és Euler közelítés esetén.

\begin{figure}
\begin{center}
\includegraphics[width=0.9\textwidth]{@EULER@}
  \caption{\textit{Euler módszer}}
\label{fig:euler}
\end{center}
\end{figure}

\subsection{Runge-Kutta 4 módszer}
\label{subsec:rk4_2}
A kiértékelés során a program egy tömbbe írta a számított adatokat, majd azokat kiírta az rk4.dat adatfileba. A tömb egyszerre tartalmazta az időt, a Föld és Hold x és y koordinátáit, illetve az x és y irányú sebességkomponenseit. A Hold koordinátáit ábrázolva látható a \ref{fig:rk4}. ábrán, hogy sokkal szebben záródik a pálya, mint az Euler módszernél.

\begin{figure}
\begin{center}
\includegraphics[width=0.9\textwidth]{@RK4@}
  \caption{\textit{Runge-Kutta 4 módszer}}
\label{fig:rk4}
\end{center}
\end{figure}

%%%%%%%%%%%%%%%%%%%%%%%%%%%%%%%%%%%%%%%%%%%%%%%%%%%%%%%%%%%%%%%%%%%%%%%%%%%%%%%%%%%%%%%%%%
% ez ugyanazt tudja, mint a \newpage, azzal a különbséggel, hogy hogyha vannak még függő 
% ábrák vagy táblázatok, amik még nem lettek elhelyezve, akkor az most itt megtörténik, és 
% csak utána megyünk tovább. A \newpage most nekünk pl hatástalan lenne.

\cleardoublepage
\vspace*{2cm}

% Mivel *-al jelöljük meg az új section-t, nem fog bekerülni a tartalomjegyzékbe, hacsak kézzel hozzá nem adjuk
\addcontentsline{toc}{section}{Köszönetnyilvánítás}
\section*{Köszönetnyilvánítás}
Szeretném megköszönni a segítséget a gyakorlatvezetőknek, Bíró Gábornak és Nagy-Egri Máténak!

%%%%%%%%%%%%%%%%%%%%%%%%%%%%%%%%%%%%%%%%%%%%%%%%%%%%%%%%%%%%%%%%%%%%%%%%%%%%%%%%%%%%%%%%%%5

\newpage \vspace*{2cm}

\addcontentsline{toc}{section}{Irodalomjegyzék}

\bibliographystyle{src/abeld}       
% van egy csomó féle/fajta stílus, google segít bennük. Én ezt használtam a diplomamunkámhoz. 
% A helyes működéshez az kell, hogy az abeld.bst ugyanebben a mappában legyen.
\bibliography{src/references}       
% a references.bib tartalmazza egyenként a hivatkozásokat, így a tex fájlod tisztább maradhat.
%
% 
%%%%%%%%%%%%%%%%%%%%%%%%%%%%%%%%%%%%%%%%%%%%%%%%%%%%%%%%%%%%%%%%%%%%%%%%%%%%%%%%%%%%%%%%

% ----------------------------------- A P P E N D I X -----------------------------------

%\newpage \vspace*{2cm}

%\pagestyle{fancy}

%\appendix
%\section{Ez egy függelék}
%\label{app:fugg}

%És igen, ezeket is lehet \ref{app:fugg}. módon hivatkozni.

% -------------------------------- N Y I L A T K O Z A T --------------------------------

% Nagyon hosszú dokumentumoknál (pl könyv) érdemes több részre bontani a forráskódot is.
% Az \input parancs után következő fájl egyszerűen "bemásolódik", így a fordító egyben 
% fogja látni az egészet.
\newpage
\input src/nyilatkozat.tex  % EZT NE HAGYD KI! OTT KELL LENNIE A DOLGOZAT VÉGÉN!


\end{document}

Ez pedig már nem fog megjelenni, mivel az end document után vagyunk.
